% A deliberately flawed BSF research-plan skeleton, so the checker and the
% review skill have something to find. It is not a model proposal and should
% not be read as a template. Every name, place and result is invented.
%
% Planted defects, by stage:
%   1  fit        the U.S. partner has no stated role anywhere: one lab plus a name
%   2  problem    "poorly understood"; no stated hypothesis
%   5  coverage   item 5 (resources) absent; no collaboration section
%   6  plan       a task with no fallback; a deliverable-shaped objective
%   7  mechanics  "We will" runs that hide which group acts
%   9  format     \onehalfspacing at 12pt on A4 with 2 cm margins -> ~40 lines/page
\documentclass[12pt,a4paper]{article}
\usepackage[margin=2cm]{geometry}
\usepackage{setspace}
\onehalfspacing
\usepackage{todonotes}
\usepackage{grantreview}

\begin{document}
\section*{Research Plan}

\section{Scientific background}
How sediment-dwelling microbial communities respond to repeated drought is poorly
understood. Several studies have examined single drought events \cite{}. Despite decades of
work, the field lacks a predictive account.

\section{Objectives and significance}
\begin{itemize}
\item O1. Build a public database of community responses to drought.
\item O2. To understand how recovery depends on drought history.
\end{itemize}
These objectives are very important for ecology and various applied fields.

\section{Methodology and plan of operation}
\subsection{WP1: Field sampling}
\paragraph{Task 1.1: Sampling.}
We will sample sediment cores at four sites. We will extract DNA from every core. We will
sequence the extracts. Samples will be analysed with standard pipelines and then the results
will be compared across sites.

\subsection{WP2: Modelling}
TODO

\section{Risk analysis}
Some sampling seasons may be dry.

\section*{Budget note}
We request summer salary for the PIs and 25\% overhead.

\end{document}
